\documentclass[../../main.tex]{subfiles}
\begin{document}
\section{Lie 群的 Lie 代数}\label{section:ch08-Lie群的Lie代数}

Lie 括号最重要的应用之一出现在 Lie 群的背景中。设 $G$ 是 Lie 群。回忆 $G$ 通过左平移在自身上光滑且传递地作用：$L_g(h)=gh$（见\cref{example:ch07-22}(c)）。$G$ 上的向量场 $X$ 称为\textbf{左不变的}，若它在所有左平移下不变，即对每个 $g\in G$，它与自身是 $L_g$-相关的。更明确地，这意味着
\begin{align}
\mathrm{d}(L_g)_{g_0}(X_{g_0})=X_{gg_0},\quad\text{对所有 }g,g_0\in G.\label{align:ch08-12}
\end{align}
由于 $L_g$ 是微分同胚，这可以简写为对每个 $g\in G$ 有 $(L_g)_*X=X$。

\begin{figure}[ht]
\centering
\includegraphics[width=0.8\textwidth]{fig-08-5.png}
\caption{定义左不变向量场}\label{figure:ch08-left-invariant}
\end{figure}

由于 $(L_g)_*(aX+bY)=a(L_g)_*X+b(L_g)_*Y$，$G$ 上所有光滑左不变向量场的集合是 $\mathfrak{X}(G)$ 的线性子空间。但它远不止于此。关键事实是它在 Lie 括号下封闭。

\begin{proposition}\label{proposition:ch08-33}
设 $G$ 是 Lie 群，$X$ 和 $Y$ 是 $G$ 上的光滑左不变向量场。则 $[X,Y]$ 也是左不变的。
\end{proposition}

\begin{proof}
设 $g\in G$ 任意。由左不变性的定义，$(L_g)_*X=X$ 且 $(L_g)_*Y=Y$，故由\cref{corollary:ch08-31}，
\begin{align*}
(L_g)_*[X,Y]=[(L_g)_*X,(L_g)_*Y]=[X,Y].
\end{align*}
因此，对每个 $g$，$[X,Y]$ 与自身是 $L_g$-相关的，即它是左不变的。
\end{proof}

一个（$\mathbb{R}$ 上的）\textbf{Lie 代数}是一个实向量空间 $\mathfrak{g}$，赋予一个称为\textbf{括号}的从 $\mathfrak{g}\times\mathfrak{g}$ 到 $\mathfrak{g}$ 的映射，通常记作 $(X,Y)\mapsto[X,Y]$，对所有 $X,Y,Z\in\mathfrak{g}$ 满足下列性质：
\begin{enumerate}[(i)]
\item \textbf{双线性}：对 $a,b\in\mathbb{R}$，
\begin{align*}
[aX+bY,Z]&=a[X,Z]+b[Y,Z],\\
[Z,aX+bY]&=a[Z,X]+b[Z,Y].
\end{align*}
\item \textbf{反称性}：
\begin{align*}
[X,Y]=-[Y,X].
\end{align*}
\item \textbf{Jacobi 恒等式}：
\begin{align*}
[X,[Y,Z]]+[Y,[Z,X]]+[Z,[X,Y]]=0.
\end{align*}
\end{enumerate}

注意 Jacobi 恒等式是结合律的替代品，后者在 Lie 代数的括号中一般不成立。在某些情形下定义 $\mathbb{C}$ 或其他域上的 Lie 代数是有用的，但我们没有理由考虑此类 Lie 代数；因此，此后我们所有的 Lie 代数都默认是实的，不再另行说明。

若 $\mathfrak{g}$ 是 Lie 代数，则线性子空间 $\mathfrak{h}\subseteq\mathfrak{g}$ 称为 $\mathfrak{g}$ 的\textbf{Lie 子代数}，若它在括号下封闭。此时 $\mathfrak{h}$ 自身在同一括号的限制下是 Lie 代数。

若 $\mathfrak{g}$ 和 $\mathfrak{h}$ 是 Lie 代数，则线性映射 $A\colon\mathfrak{g}\to\mathfrak{h}$ 称为\textbf{Lie 代数同态}，若它保持括号：$A[X,Y]=[AX,AY]$。可逆的 Lie 代数同态称为\textbf{Lie 代数同构}。若存在从 $\mathfrak{g}$ 到 $\mathfrak{h}$ 的 Lie 代数同构，则称它们作为 Lie 代数是同构的。

\begin{exercise}\label{exercise:ch08-34}
验证 Lie 代数同态的核和像是 Lie 子代数。
\end{exercise}

\begin{exercise}\label{exercise:ch08-35}
假设 $\mathfrak{g}$ 和 $\mathfrak{h}$ 是有限维 Lie 代数，$A\colon\mathfrak{g}\to\mathfrak{h}$ 是线性映射。证明 $A$ 是 Lie 代数同态，当且仅当对 $\mathfrak{g}$ 的某组基 $(E_1,\dots,E_n)$ 有 $A[E_i,E_j]=[AE_i,AE_j]$。
\end{exercise}

\begin{example}\label{example:ch08-36}[Lie 代数]
\begin{enumerate}[(a)]
\item 光滑流形 $M$ 上所有光滑向量场的空间 $\mathfrak{X}(M)$ 在 Lie 括号下由\cref{proposition:ch08-28}可知是 Lie 代数。

\item 若 $G$ 是 Lie 群，则 $G$ 上所有光滑左不变向量场的集合是 $\mathfrak{X}(G)$ 的 Lie 子代数，因而是 Lie 代数。

\item $n\times n$ 实矩阵的向量空间 $\mathcal{M}(n,\mathbb{R})$ 在交换子括号
\begin{align*}
[A,B]=AB-BA
\end{align*}
下成为 $n^2$ 维 Lie 代数。双线性和反称性由定义显然，Jacobi 恒等式由直接计算可得。当我们将 $\mathcal{M}(n,\mathbb{R})$ 视为带有此括号的 Lie 代数时，记作 $\mathfrak{gl}(n,\mathbb{R})$。

\item 类似地，$\mathfrak{gl}(n,\mathbb{C})$ 是在 $\mathcal{M}(n,\mathbb{C})$ 上赋予交换子括号得到的 $2n^2$ 维（实）Lie 代数。

\item 若 $V$ 是向量空间，则从 $V$ 到自身的所有线性映射的向量空间在交换子括号
\begin{align*}
[A,B]=A\circ B-B\circ A
\end{align*}
下成为 Lie 代数，记作 $\mathfrak{gl}(V)$。在我们通常将 $n\times n$ 矩阵与从 $\mathbb{R}^n$ 到自身的线性映射等同下，$\mathfrak{gl}(\mathbb{R}^n)$ 与 $\mathfrak{gl}(n,\mathbb{R})$ 相同。

\item 任意向量空间 $V$ 若定义所有括号为零，则成为 Lie 代数。此类 Lie 代数称为\textbf{交换的}。（该名称反映了这样一个事实：正如前面的例子，大多数 Lie 代数中的括号是由底层结合积定义的交换子，因此所有括号恰好在底层积交换时为零；它也反映了交换 Lie 代数与交换 Lie 群之间的联系，你将在\cref{exercise:ch08-25}和\cref{exercise:ch20-07}中探讨。）
\end{enumerate}
\end{example}

\cref{example:ch08-36}(b) 是最重要的一个。Lie 群 $G$ 上所有光滑左不变向量场的 Lie 代数称为 $G$ 的\textbf{Lie 代数}，记作 $\operatorname{Lie}(G)$。（下面将看到光滑性假设是多余的；见\cref{corollary:ch08-38}。）基本事实是 $\operatorname{Lie}(G)$ 是有限维的，事实上其维数与 $G$ 本身相同，如下定理所示。

\begin{theorem}\label{theorem:ch08-37}
设 $G$ 是 Lie 群。赋值映射 $\varepsilon\colon\operatorname{Lie}(G)\to T_eG$，由 $\varepsilon(X)=X_e$ 给出，是向量空间同构。因此，$\operatorname{Lie}(G)$ 是有限维的，维数等于 $\dim G$。
\end{theorem}

\begin{proof}
由定义显然 $\varepsilon$ 是 $\mathbb{R}$-线性的。容易证明它是单射：若对某个 $X\in\operatorname{Lie}(G)$ 有 $\varepsilon(X)=X_e=0$，则 $X$ 的左不变性蕴涵对每个 $g\in G$ 有 $X_g=\mathrm{d}(L_g)_e(X_e)=0$，故 $X=0$。

为证明 $\varepsilon$ 是满射，任取 $v\in T_eG$，在 $G$ 上定义一个（粗）向量场 $v^L$ 为
\begin{align}
(v^L)_g=\mathrm{d}(L_g)_e(v).\label{align:ch08-13}
\end{align}
若 $G$ 上存在在单位元处取值为 $v$ 的左不变向量场，显然它必须由该公式给出。

首先需要验证 $v^L$ 是光滑的。由\cref{proposition:ch08-14}，只需证明每当 $f\in C^\infty(G)$ 时 $v^Lf$ 光滑。选取光滑曲线 $\gamma\colon(-\delta,\delta)\to G$ 使得 $\gamma(0)=e$ 且 $\gamma'(0)=v$。则对所有 $g\in G$，
\begin{align*}
(v^Lf)(g)=(v^L)_gf=\mathrm{d}(L_g)_e(v)f=v(f\circ L_g)=\gamma'(0)(f\circ L_g)=\left.\frac{\mathrm{d}}{\mathrm{d}t}\right|_{t=0}(f\circ L_g\circ\gamma)(t).
\end{align*}
若定义 $\Psi\colon(-\delta,\delta)\times G\to\mathbb{R}$ 为 $\Psi(t,g)=f\circ L_g\circ\gamma(t)=f(g\gamma(t))$，上述计算表明 $(v^Lf)(g)=\partial\Psi/\partial t(0,g)$。由于 $\Psi$ 是群乘法、$f$ 和 $\gamma$ 的复合，它是光滑的。由此可知 $\partial\Psi/\partial t(0,g)$ 光滑地依赖于 $g$，故 $v^Lf$ 光滑。

接下来证明 $v^L$ 是左不变的，即对所有 $g,h\in G$ 有 $\mathrm{d}(L_h)_g((v^L)_g)=(v^L)_{hg}$。这由 $v^L$ 的定义以及 $L_h\circ L_g=L_{hg}$ 可得：
\begin{align*}
\mathrm{d}(L_h)_g\bigl((v^L)_g\bigr)=\mathrm{d}(L_h)_g\circ\mathrm{d}(L_g)_e(v)=\mathrm{d}(L_h\circ L_g)_e(v)=\mathrm{d}(L_{hg})_e(v)=(v^L)_{hg}.
\end{align*}
因此 $v^L\in\operatorname{Lie}(G)$。由于 $L_e$（由单位元的左平移）是 $G$ 的恒等映射，故 $\varepsilon(v^L)=v^L|_e=v$，从而 $\varepsilon$ 是满射。
\end{proof}

给定任意向量 $v\in T_eG$，我们继续使用记号 $v^L$ 表示由\eqref{align:ch08-13}定义的光滑左不变向量场。

值得注意的是，上述证明也表明 $\operatorname{Lie}(G)$ 定义中的光滑性假设是不必要的。

\begin{corollary}\label{corollary:ch08-38}
Lie 群上的每个左不变粗向量场都是光滑的。
\end{corollary}

\begin{proof}
设 $X$ 是 Lie 群 $G$ 上的左不变粗向量场，令 $v=X_e$。$X$ 左不变这一事实蕴涵 $X=v^L$，而后者是光滑的。
\end{proof}

全局左不变向量场的存在性还给出 Lie 群如下重要性质。回忆，光滑流形若容许光滑全局标架则称为可平行化的。若 $G$ 是 Lie 群，则由左不变向量场组成的局部或全局标架称为\textbf{左不变标架}。

\begin{corollary}\label{corollary:ch08-39}
每个 Lie 群都容许左不变的光滑全局标架，因此每个 Lie 群都是可平行化的。
\end{corollary}

\begin{proof}
若 $G$ 是 Lie 群，则 $\operatorname{Lie}(G)$ 的每组基都是 $G$ 的左不变光滑全局标架。
\end{proof}

\begin{example}\label{example:ch08-40}
我们来确定一些熟悉 Lie 群的 Lie 代数。
\begin{enumerate}[(a)]
\item \textbf{Euclid 空间 $\mathbb{R}^n$}：若将 $\mathbb{R}^n$ 视为加法下的 Lie 群，则由元素 $b\in\mathbb{R}^n$ 的左平移由仿射映射 $L_b(x)=b+x$ 给出，其微分 $\mathrm{d}(L_b)$ 在标准坐标下由单位矩阵表示。因此，向量场 $X^i\partial/\partial x^i$ 是左不变的，当且仅当其系数 $X^i$ 为常数。由于由\eqref{align:ch08-08}可知两个常系数向量场的 Lie 括号为零，$\mathbb{R}^n$ 的 Lie 代数是交换的，同构于带平凡括号的 $\mathbb{R}^n$ 本身。简而言之，$\operatorname{Lie}(\mathbb{R}^n)\cong\mathbb{R}^n$。

\item \textbf{圆周群 $S^1$}：在适当的角坐标下，每个左平移有形如 $\theta\mapsto\theta+c$ 的局部坐标表示。由于该映射的微分是 $1\times1$ 单位矩阵，由此可知\cref{example:ch08-04}中定义的向量场 $\mathrm{d}/\mathrm{d}\theta$ 是左不变的，因而是 $S^1$ 的 Lie 代数的一组基。该 Lie 代数是一维交换的，因此 $\operatorname{Lie}(S^1)\cong\mathbb{R}$。

\item \textbf{$n$ 维环面 $\mathbb{T}^n=S^1\times\cdots\times S^1$}：与上述类似的分析表明 $(\partial/\partial\theta_1,\dots,\partial/\partial\theta_n)$ 是 $\operatorname{Lie}(\mathbb{T}^n)$ 的一组基，其中 $\partial/\partial\theta_i$ 是第 $i$ 个 $S^1$ 因子上的角坐标向量场。由于这些坐标向量场的 Lie 括号全为零，$\operatorname{Lie}(\mathbb{T}^n)\cong\mathbb{R}^n$。
\end{enumerate}
\end{example}

Lie 群 $\mathbb{R}^n$、$S^1$ 和 $\mathbb{T}^n$ 是交换的，且如上面的讨论所示，它们的 Lie 代数也是交换的。这并非偶然：每个交换 Lie 群都有交换的 Lie 代数（见\cref{exercise:ch08-25}）。稍后将看到，反过来在群连通的条件下也成立（\cref{exercise:ch20-07}）。

正如我们可以将切空间视为光滑流形在一点附近的"线性模型"，Lie 群的 Lie 代数提供了群的"线性模型"，它反映群的许多性质。由于 Lie 群比普通光滑流形具有更多结构，其线性模型比普通向量空间具有更多结构也就不足为奇了。由于有限维 Lie 代数是纯线性代数对象，它在许多方面比群本身更易于理解。Lie 群理论的许多进展都来自对 Lie 代数的仔细分析。

本节最后分析最重要的非交换 Lie 群——一般线性群——的 Lie 代数。\cref{theorem:ch08-37}给出了 $\operatorname{Lie}(\operatorname{GL}(n,\mathbb{R}))$ 与 $\operatorname{GL}(n,\mathbb{R})$ 在单位矩阵 $I_n$ 处的切空间之间的向量空间同构。由于 $\operatorname{GL}(n,\mathbb{R})$ 是向量空间 $\mathfrak{gl}(n,\mathbb{R})$ 的开子集，其切空间自然同构于 $\mathfrak{gl}(n,\mathbb{R})$ 本身。这两个同构的复合给出向量空间同构 $\operatorname{Lie}(\operatorname{GL}(n,\mathbb{R}))\cong\mathfrak{gl}(n,\mathbb{R})$。

向量空间 $\operatorname{Lie}(\operatorname{GL}(n,\mathbb{R}))$ 和 $\mathfrak{gl}(n,\mathbb{R})$ 各自具有独立定义的 Lie 代数结构——前者来自向量场的 Lie 括号，后者来自矩阵的交换子括号。下一个命题表明这些空间之间的自然向量空间同构实际上是 Lie 代数同构。

\begin{proposition}\label{proposition:ch08-41}[一般线性群的 Lie 代数]
自然映射的复合
\begin{align}
\operatorname{Lie}\bigl(\operatorname{GL}(n,\mathbb{R})\bigr)\to T_{I_n}\operatorname{GL}(n,\mathbb{R})\to\mathfrak{gl}(n,\mathbb{R})\label{align:ch08-14}
\end{align}
给出 $\operatorname{Lie}(\operatorname{GL}(n,\mathbb{R}))$ 与矩阵代数 $\mathfrak{gl}(n,\mathbb{R})$ 之间的 Lie 代数同构。
\end{proposition}

\begin{proof}
使用矩阵元 $X^i_j$ 作为 $\operatorname{GL}(n,\mathbb{R})\subseteq\mathfrak{gl}(n,\mathbb{R})$ 上的整体坐标，自然同构 $T_{I_n}\operatorname{GL}(n,\mathbb{R})\to\mathfrak{gl}(n,\mathbb{R})$ 形如
\begin{align*}
A^i_j\frac{\partial}{\partial X^i_j}\bigg|_{I_n}\mapsto\bigl(A^i_j\bigr).
\end{align*}
（由于指标 $i,j$ 同时作为坐标指标和矩阵行列指标，此处不可能保持我们所有坐标取上指标的约定。然而，我们继续遵守 Einstein 求和约定及其他相关指标约定。特别地，在上面的表达式中，"分母中"的上指标应视为下指标，反之亦然。）

记 $\mathfrak{g}$ 为 $\operatorname{GL}(n,\mathbb{R})$ 的 Lie 代数。任意矩阵 $A=(A^i_j)\in\mathfrak{gl}(n,\mathbb{R})$ 确定了一个由\eqref{align:ch08-13}定义的左不变向量场 $A^L\in\mathfrak{g}$，在此情形下成为
\begin{align*}
(A^L)_X=\mathrm{d}(L_X)_{I_n}(A)=\mathrm{d}(L_X)_{I_n}\left(A^i_j\frac{\partial}{\partial X^i_j}\bigg|_{I_n}\right).
\end{align*}
由于 $L_X$ 是 $\mathfrak{gl}(n,\mathbb{R})$ 上线性映射 $A\mapsto XA$ 在 $\operatorname{GL}(n,\mathbb{R})$ 上的限制，其微分在坐标中由完全相同的线性映射表示。换言之，由 $A$ 确定的左不变向量场 $A^L$ 在 $X\in\operatorname{GL}(n,\mathbb{R})$ 处的值为
\begin{align}
(A^L)_X=X^i_jA^j_k\frac{\partial}{\partial X^i_k}\bigg|_X.\label{align:ch08-15}
\end{align}

给定两个矩阵 $A,B\in\mathfrak{gl}(n,\mathbb{R})$，相应左不变向量场的 Lie 括号由
\begin{align*}
[A^L,B^L]&=\left[X^i_jA^j_k\frac{\partial}{\partial X^i_k},\,X^p_qB^q_r\frac{\partial}{\partial X^p_r}\right]\\
&=X^i_jA^j_k\frac{\partial}{\partial X^i_k}\left(X^p_qB^q_r\right)\frac{\partial}{\partial X^p_r}-X^p_qB^q_r\frac{\partial}{\partial X^p_r}\left(X^i_jA^j_k\right)\frac{\partial}{\partial X^i_k}\\
&=X^i_jA^j_kB^k_r\frac{\partial}{\partial X^i_r}-X^p_qB^q_rA^r_k\frac{\partial}{\partial X^p_k}\\
&=\left(X^i_jA^j_kB^k_r-X^i_jB^j_kA^k_r\right)\frac{\partial}{\partial X^i_r}
\end{align*}
给出，其中使用了 $\partial X^p_q/\partial X^i_k$ 在 $p=i$ 且 $q=k$ 时等于 $1$，否则等于 $0$，以及 $A^i_j$ 和 $B^i_j$ 为常数这一事实。在 $X$ 等于单位矩阵时计算最后这个表达式，得到
\begin{align*}
[A^L,B^L]_{I_n}=\left(A^i_kB^k_r-B^i_kA^k_r\right)\frac{\partial}{\partial X^i_r}\bigg|_{I_n}.
\end{align*}
这是对应于矩阵交换子括号 $[A,B]$ 的向量。由于左不变向量场 $[A^L,B^L]$ 由其在单位元处的值确定，这蕴涵
\begin{align*}
[A^L,B^L]=[A,B]^L,
\end{align*}
这正是复合映射\eqref{align:ch08-14}是 Lie 代数同构的陈述。
\end{proof}

对抽象向量空间有该结果的类似物。若 $V$ 是任意有限维实向量空间，回忆我们已将 $\operatorname{GL}(V)$ 定义为 $V$ 的可逆线性变换的 Lie 群，将 $\mathfrak{gl}(V)$ 定义为所有线性变换的 Lie 代数。正如 $\operatorname{GL}(n,\mathbb{R})$ 的情形，可将 $\operatorname{GL}(V)$ 视为 $\mathfrak{gl}(V)$ 的开子流形，因此存在典范向量空间同构
\begin{align}
\operatorname{Lie}\bigl(\operatorname{GL}(V)\bigr)\to T_{\operatorname{Id}}\operatorname{GL}(V)\to\mathfrak{gl}(V).\label{align:ch08-16}
\end{align}

\begin{corollary}\label{corollary:ch08-42}
若 $V$ 是任意有限维实向量空间，则\eqref{align:ch08-16}中典范同构的复合给出 $\operatorname{Lie}(\operatorname{GL}(V))$ 与 $\mathfrak{gl}(V)$ 之间的 Lie 代数同构。
\end{corollary}

\begin{exercise}\label{exercise:ch08-43}
通过为 $V$ 选取一组基并应用\cref{proposition:ch08-41}来证明前述推论。
\end{exercise}

\textbf{诱导 Lie 代数同态}

Lie 群 Lie 代数的重要性很大程度上源于如下事实：每个 Lie 群同态都诱导一个 Lie 代数同态，如下一个定理所示。

\begin{theorem}\label{theorem:ch08-44}[诱导 Lie 代数同态]
设 $G$ 和 $H$ 是 Lie 群，$\mathfrak{g}$ 和 $\mathfrak{h}$ 是它们的 Lie 代数。假设 $F\colon G\to H$ 是 Lie 群同态。对每个 $X\in\mathfrak{g}$，存在 $\mathfrak{h}$ 中唯一的向量场与 $X$ 是 $F$-相关的。将该向量场记作 $F_*X$，如此定义的映射 $F_*\colon\mathfrak{g}\to\mathfrak{h}$ 是 Lie 代数同态。
\end{theorem}

\begin{figure}[ht]
\centering
\includegraphics[width=0.8\textwidth]{fig-08-6.png}
\caption{诱导的 Lie 代数同态}\label{figure:ch08-induced-lie-algebra}
\end{figure}

\begin{proof}
若存在任何与 $X$ 是 $F$-相关的向量场 $Y\in\mathfrak{h}$，它必须满足 $Y_e=\mathrm{d}F_e(X_e)$，因而必须由
\begin{align*}
Y=(\mathrm{d}F_e(X_e))^L
\end{align*}
唯一确定。

为证明此 $Y$ 与 $X$ 是 $F$-相关的，注意 $F$ 是同态这一事实蕴涵
\begin{align*}
F(gg')=F(g)F(g')&\Rightarrow F(L_gg')=L_{F(g)}F(g')\\
&\Rightarrow F\circ L_g=L_{F(g)}\circ F\\
&\Rightarrow \mathrm{d}F\circ\mathrm{d}(L_g)=\mathrm{d}(L_{F(g)})\circ\mathrm{d}F.
\end{align*}
因此，
\begin{align*}
\mathrm{d}F(X_g)&=\mathrm{d}F\bigl(\mathrm{d}(L_g)(X_e)\bigr)\\
&=\mathrm{d}\bigl(L_{F(g)}\bigr)\bigl(\mathrm{d}F(X_e)\bigr)\\
&=\mathrm{d}\bigl(L_{F(g)}\bigr)(Y_e)=Y_{F(g)}.
\end{align*}
这恰好说明 $X$ 和 $Y$ 是 $F$-相关的。

对每个 $X\in\mathfrak{g}$，令 $F_*X$ 表示 $\mathfrak{h}$ 中与 $X$ 是 $F$-相关的唯一向量场。于是由 Lie 括号的自然性立即可得 $F_*[X,Y]=[F_*X,F_*Y]$，故 $F_*$ 是 Lie 代数同态。
\end{proof}

该定理所断言存在的映射 $F_*\colon\mathfrak{g}\to\mathfrak{h}$ 称为\textbf{诱导 Lie 代数同态}。注意该定理蕴涵对任意左不变向量场 $X\in\mathfrak{g}$，$F_*X$ 是 $H$ 上良定义的光滑向量场，即使 $F$ 可能不是微分同胚。

\begin{proposition}\label{proposition:ch08-45}[诱导同态的性质]
\begin{enumerate}[(a)]
\item 由 $G$ 的恒等映射诱导的同态 $(\operatorname{Id}_G)_*\colon\operatorname{Lie}(G)\to\operatorname{Lie}(G)$ 是 $\operatorname{Lie}(G)$ 的恒等映射。
\item 若 $F_1\colon G\to H$ 和 $F_2\colon H\to K$ 是 Lie 群同态，则
\begin{align*}
(F_2\circ F_1)_*=(F_2)_*\circ(F_1)_*\colon\operatorname{Lie}(G)\to\operatorname{Lie}(K).
\end{align*}
\item 同构的 Lie 群有同构的 Lie 代数。
\end{enumerate}
\end{proposition}

\begin{proof}
微分同时满足 $\mathrm{d}(\operatorname{Id}_G)_e=\operatorname{Id}_{T_eG}$ 和 $\mathrm{d}(F_2\circ F_1)_e=\mathrm{d}(F_2)_e\circ\mathrm{d}(F_1)_e$。由于诱导同态由单位元处的微分确定，这证明了 (a) 和 (b)。若 $F\colon G\to H$ 是同构，则 (a) 和 (b) 一起蕴涵 $F_*\circ(F^{-1})_*=(F\circ F^{-1})_*=\operatorname{Id}=(F^{-1})_*\circ F_*$，故 $F_*\colon\operatorname{Lie}(G)\to\operatorname{Lie}(H)$ 是同构。
\end{proof}

你可能已经猜到，前述命题有一种范畴解释。记 $\mathbf{Lie}$ 为以 Lie 群为对象、Lie 群同态为态射的范畴，$\mathbf{lie}$ 为以有限维 Lie 代数为对象、Lie 代数同态为态射的范畴。\cref{proposition:ch08-45}可解释为表明指派 $G\mapsto\operatorname{Lie}(G)$、$F\mapsto F_*$ 定义了从 $\mathbf{Lie}$ 到 $\mathbf{lie}$ 的共变函子。

\textbf{Lie 子群的 Lie 代数}

若 $G$ 是 Lie 群，$H\subseteq G$ 是 Lie 子群，我们可能希望 $H$ 的 Lie 代数是 $G$ 的 Lie 代数的 Lie 子代数。然而，$\operatorname{Lie}(H)$ 的元素是 $H$ 上的向量场，而非 $G$ 上的，因此严格来说不是 $\operatorname{Lie}(G)$ 的元素。尽管如此，下一个命题给出了将 $\operatorname{Lie}(H)$ 视为 $\operatorname{Lie}(G)$ 的子代数的方法。

\begin{theorem}\label{theorem:ch08-46}[Lie 子群的 Lie 代数]
假设 $H\subseteq G$ 是 Lie 子群，$\iota\colon H\hookrightarrow G$ 是包含映射。存在 $\operatorname{Lie}(G)$ 的 Lie 子代数 $\mathfrak{h}$，它典范同构于 $\operatorname{Lie}(H)$，由下列任一特征刻画：
\begin{align}
\mathfrak{h}=\iota_*\bigl(\operatorname{Lie}(H)\bigr)=\{X\in\operatorname{Lie}(G):X_e\in T_eH\}.\label{align:ch08-17}
\end{align}
\end{theorem}

\begin{proof}
由于包含映射 $\iota\colon H\hookrightarrow G$ 是 Lie 群同态，$\iota_*(\operatorname{Lie}(H))$ 是 $\operatorname{Lie}(G)$ 的 Lie 子代数。由我们定义诱导 Lie 代数同态的方式，该子代数恰为 $G$ 上在单位元处的值形如 $\mathrm{d}\iota_e(v)$（其中 $v\in T_eH$）的左不变向量场的集合。由于微分 $\mathrm{d}\iota_e\colon T_eH\to T_eG$ 是将 $T_eH$ 作为子空间包含入 $T_eG$，\eqref{align:ch08-17}中 $\mathfrak{h}$ 的两个特征刻画相等。由于 $\mathrm{d}\iota_e$ 在 $T_eH$ 上单射，由此可知 $\iota_*$ 在 $\operatorname{Lie}(H)$ 上单射；又由于由 $\mathfrak{h}$ 的定义它是满射，故它是 $\operatorname{Lie}(H)$ 与 $\mathfrak{h}$ 之间的同构。
\end{proof}

使用该命题，每当 $H$ 是 $G$ 的 Lie 子群时，我们常将 $\operatorname{Lie}(H)$ 视为 $\operatorname{Lie}(G)$ 的子代数。如上面提到的，$\operatorname{Lie}(H)$ 的元素本身不是 $G$ 上的左不变向量场。但前述命题表明，$\operatorname{Lie}(H)$ 的每个元素对应于 $\operatorname{Lie}(G)$ 中由其在单位元处的值确定的唯一元素，而如此确定的从 $\operatorname{Lie}(H)$ 到 $\operatorname{Lie}(G)$ 的单射保持 Lie 括号；因此将 $\operatorname{Lie}(H)$ 视为 $\operatorname{Lie}(G)$ 的子代数并无大错。

在 $\operatorname{GL}(n,\mathbb{R})$ 的 Lie 子群情形中，这种等同尤其有启发性。

\begin{example}\label{example:ch08-47}[$O(n)$ 的 Lie 代数]
正交群 $O(n)$ 是 $\operatorname{GL}(n,\mathbb{R})$ 的 Lie 子群。由\cref{example:ch07-27}，它等于水平集 $\Phi^{-1}(I_n)$，其中 $\Phi\colon\operatorname{GL}(n,\mathbb{R})\to\mathcal{M}(n,\mathbb{R})$ 是映射 $\Phi(A)=A^{T}A$。由\cref{exercise:ch05-40}的结果，$T_{I_n}O(n)$ 等于微分 $\mathrm{d}\Phi_{I_n}\colon T_{I_n}\operatorname{GL}(n,\mathbb{R})\to T_{I_n}\mathcal{M}(n,\mathbb{R})$ 的核。由\cref{example:ch07-27}中的计算，该微分为 $\mathrm{d}\Phi_{I_n}(B)=B^{T}+B$，故
\begin{align*}
T_{I_n}O(n)=\{B\in\mathfrak{gl}(n,\mathbb{R}):B^{T}+B=0\}=\{\text{反称 }n\times n\text{ 矩阵}\}.
\end{align*}
我们将 $\mathfrak{gl}(n,\mathbb{R})$ 的这个子空间记作 $\mathfrak{o}(n)$。\cref{theorem:ch08-46}表明 $\mathfrak{o}(n)$ 是 $\mathfrak{gl}(n,\mathbb{R})$ 的 Lie 子代数，典范同构于 $\operatorname{Lie}(O(n))$。注意我们甚至不必直接验证 $\mathfrak{o}(n)$ 是子代数。
\end{example}

上面证明了 $\operatorname{GL}(n,\mathbb{R})$ 的 Lie 代数自然同构于矩阵代数 $\mathfrak{gl}(n,\mathbb{R})$。现在可以对 $\operatorname{GL}(n,\mathbb{C})$ 证明类似结果。正如实情形，我们通常将 $\operatorname{GL}(n,\mathbb{C})$ 视为 $\mathfrak{gl}(n,\mathbb{C})$ 的开子集，这给出一系列向量空间同构
\begin{align}
\operatorname{Lie}\bigl(\operatorname{GL}(n,\mathbb{C})\bigr)\xrightarrow{\varepsilon}T_{I_n}\operatorname{GL}(n,\mathbb{C})\xrightarrow{\varphi}\mathfrak{gl}(n,\mathbb{C}),\label{align:ch08-18}
\end{align}
其中 $\varepsilon$ 是赋值映射，$\varphi$ 是向量空间开子集的切空间与向量空间本身之间的通常等同。（注意我们将这些视为实向量空间，而非复向量空间。）

\begin{proposition}\label{proposition:ch08-48}[$\operatorname{GL}(n,\mathbb{C})$ 的 Lie 代数]
\eqref{align:ch08-18}中映射的复合给出 $\operatorname{Lie}(\operatorname{GL}(n,\mathbb{C}))$ 与矩阵代数 $\mathfrak{gl}(n,\mathbb{C})$ 之间的 Lie 代数同构。
\end{proposition}

\begin{proof}
\cref{example:ch07-18}(d)中构造的 Lie 群同态 $\beta\colon\operatorname{GL}(n,\mathbb{C})\to\operatorname{GL}(2n,\mathbb{R})$ 诱导 Lie 代数同态
\begin{align*}
\beta_*\colon\operatorname{Lie}\bigl(\operatorname{GL}(n,\mathbb{C})\bigr)\to\operatorname{Lie}\bigl(\operatorname{GL}(2n,\mathbb{R})\bigr).
\end{align*}
将 $\beta_*$ 与我们的典范同构复合，得到交换图
\begin{align*}
\begin{tikzcd}[ampersand replacement=\&]
\operatorname{Lie}(\operatorname{GL}(n,\mathbb{C})) \ar{r}{\varepsilon} \ar{d}{\beta_*} \& T_{I_n}\operatorname{GL}(n,\mathbb{C}) \ar{r}{\varphi} \ar{d}{\mathrm{d}\beta_{I_n}} \& \mathfrak{gl}(n,\mathbb{C}) \ar{d}{\alpha}\\
\operatorname{Lie}(\operatorname{GL}(2n,\mathbb{R})) \ar{r}{\varepsilon} \& T_{I_{2n}}\operatorname{GL}(2n,\mathbb{R}) \ar{r}{\varphi} \& \mathfrak{gl}(2n,\mathbb{R})
\end{tikzcd}
\end{align*}
其中 $\alpha=\varphi\circ\mathrm{d}\beta_{I_n}\circ\varphi^{-1}$。\cref{proposition:ch08-41}表明底行同构的复合是 Lie 代数同构；我们需要对顶行证明同样的结论。

由\cref{example:ch07-18}(d)中的公式容易看出 $\beta$ 是（线性映射的）限制。由此可知 $\mathrm{d}\beta_{I_n}\colon T_{I_n}\operatorname{GL}(n,\mathbb{C})\to T_{I_{2n}}\operatorname{GL}(2n,\mathbb{R})$ 与 $\beta$ 由完全相同的公式给出，$\alpha\colon\mathfrak{gl}(n,\mathbb{C})\to\mathfrak{gl}(2n,\mathbb{R})$ 亦然。由于 $\beta(AB)=\beta(A)\beta(B)$，由此可知 $\alpha$ 保持矩阵交换子：
\begin{align*}
\alpha[A,B]=\alpha(AB-BA)=\alpha(A)\alpha(B)-\alpha(B)\alpha(A)=[\alpha(A),\alpha(B)].
\end{align*}
因此 $\alpha$ 是从 $\mathfrak{gl}(n,\mathbb{C})$ 到 $\mathfrak{gl}(2n,\mathbb{R})$ 的单射 Lie 代数同态（将两者视为矩阵代数）。将\eqref{align:ch08-19}中的底行替换为竖直映射的像，得到向量空间同构的交换图
\begin{align*}
\begin{tikzcd}[ampersand replacement=\&]
\operatorname{Lie}(\operatorname{GL}(n,\mathbb{C})) \ar{r}{\cong} \ar{d}{\beta_*} \& \mathfrak{gl}(n,\mathbb{C}) \ar{d}{\alpha}\\
\beta_*(\operatorname{Lie}(\operatorname{GL}(n,\mathbb{C}))) \ar{r}{\cong} \& \alpha(\mathfrak{gl}(n,\mathbb{C}))
\end{tikzcd}
\end{align*}
其中底行映射和两个竖直映射都是 Lie 代数同构；由此可知顶行映射也是 Lie 代数同构。
\end{proof}

与 Lie 群表示的概念平行，也有 Lie 代数表示的概念。若 $\mathfrak{g}$ 是有限维 Lie 代数，则 $\mathfrak{g}$ 的（有限维）\textbf{表示}是从 $\mathfrak{g}$ 到某个有限维向量空间 $V$ 的 $\mathfrak{gl}(V)$ 的 Lie 代数同态 $\varphi\colon\mathfrak{g}\to\mathfrak{gl}(V)$，其中 $\mathfrak{gl}(V)$ 表示从 $V$ 到自身的线性映射的 Lie 代数。若 $\varphi$ 是单射的，则称它是\textbf{忠实的}，此时 $\mathfrak{g}$ 同构于 Lie 子代数 $\varphi(\mathfrak{g})\subseteq\mathfrak{gl}(V)\cong\mathfrak{gl}(n,\mathbb{R})$。

Lie 群表示与其 Lie 代数表示之间有密切联系。假设 $G$ 是 Lie 群，$\mathfrak{g}$ 是其 Lie 代数。若 $\rho\colon G\to\operatorname{GL}(V)$ 是 $G$ 的任意表示，则容易验证 $\rho_*\colon\mathfrak{g}\to\mathfrak{gl}(V)$ 是 $\mathfrak{g}$ 的表示。

本节以陈述一个关于 Lie 代数的深刻代数结果作结，我们将在第 20 章中使用它。下列定理的证明需要远超我们所掌握的代数知识，因此建议感兴趣的读者参阅 [Var84] 中的证明。

\begin{theorem}\label{theorem:ch08-49}[Ado 定理]
每个有限维实 Lie 代数都容许忠实的有限维表示。
\end{theorem}

\begin{corollary}\label{corollary:ch08-50}
每个有限维实 Lie 代数都同构于某个带交换子括号的矩阵代数 $\mathfrak{gl}(n,\mathbb{R})$ 的 Lie 子代数。
\end{corollary}

\begin{proof}
设 $\mathfrak{g}$ 是有限维实 Lie 代数。由 Ado 定理，$\mathfrak{g}$ 对某个有限维实向量空间 $V$ 有忠实表示 $\rho\colon\mathfrak{g}\to\mathfrak{gl}(V)$。为 $V$ 选取一组基便给出 $\mathfrak{gl}(V)$ 与某个 $n$ 的 $\mathfrak{gl}(n,\mathbb{R})$ 之间的同构，将 $\rho$ 与此同构复合得到从 $\mathfrak{g}$ 到 $\mathfrak{gl}(n,\mathbb{R})$ 的单射 Lie 代数同态。其像是同构于 $\mathfrak{g}$ 的 Lie 子代数。
\end{proof}

如前一章所提到的，重要的是要记住 Lie 群的类似结果是错误的：存在不同构于 $\operatorname{GL}(n,\mathbb{R})$ 的 Lie 子群的 Lie 群（例子见\cref{exercise:ch21-26}）。

\end{document}